\section{Related Work}
Our work extends FVAPPS by Dougherty and Mehta \cite{dougherty2025proving}, which pioneered AI-driven translation of coding problems into Lean theorems. In FVAPPS, source programs were taken from the APPS benchmark \cite{hendrycks2021apps}, translated to PBTs, then lifted to theorems. We apply the same process at scale to PBTs found in the wild, rather than synthetic coding challenges.

\paragraph{Lean verification benchmarks.}
A growing number of benchmarks target Lean specifically. miniCodeProps \cite{brandfonbrener2024minicodeprops} provides 177 program specifications translated from TIP (Tons of Inductive Programs), finding that current LLM provers fail on medium and hard properties. CLEVER \cite{materzok2025clever} offers 161 curated problems sourced from HumanEval, testing end-to-end verified code generation. Verina \cite{wu2025verina} covers all three foundational tasks---code generation, specification generation, and proof generation---across 189 manually curated coding tasks. VeriBench \cite{stechly2025veribench} translates Python code with documentation into Lean implementations, specifications, and unit tests (113 tasks). The Vericoding Benchmark \cite{bursuc2025benchmarkvericodingformallyverified} supports Lean, Rust/Verus, and Dafny, reporting success rates from 27\% (Lean) to 82\% (Dafny). VeriEquivBench \cite{ma2025veriequivbench} uses LeetCode transformations to evaluate verifiable agents on novel data without contamination. These benchmarks are valuable but share a common limitation: their specifications derive from synthetic coding problems or pedagogical exercises, not from properties written by practicing engineers.

\paragraph{Other verification tool benchmarks.}
DafnyBench \cite{loughridge2024dafnybench} provides Dafny verification tasks. VerusBench \cite{chen2024verusbench} collects 150 proof tasks spanning Dafny, SV-COMP, and Verus. SV-COMP \cite{beyer2025svcomp} is the standard competition for software verifiers, with over 33{,}000 C tasks, but targets automated verification rather than interactive theorem proving. InvBench \cite{wu2025invbench} focuses specifically on loop invariant synthesis. Chakraborty et al.\ \cite{chakraborty2024neural} explore neural synthesis for SMT-assisted proof-oriented programming in F*.

\paragraph{Interactive theorem proving benchmarks.}
LeanDojo Benchmark 4 \cite{yang2024leandojo} provides 122{,}517 theorems from Mathlib4. CoqStoq \cite{thompson2025rangoadaptiveretrievalaugmentedproving} collects 196{,}929 theorems from 2{,}226 GitHub Coq repositories. These are primarily mathematical rather than software-oriented, and their theorems exist in the training data of most frontier models.

\paragraph{Mathematics benchmarks.}
FrontierMath \cite{glazer2024frontiermath} evaluates advanced mathematical reasoning. While formal verification and mathematics share proof infrastructure, the skills required differ substantially: software verification demands modelling of program semantics, library APIs, and computational behavior, rather than abstract mathematical reasoning.

\FVSpec differs from existing benchmarks in three ways. First, our specifications originate from real-world property-based tests written by practicing engineers, not synthetic problems or pedagogical exercises. Second, the source PBTs were never formally verified, so most theorems in our benchmark will not appear in any model's training data. Third, we operate at a larger scale than existing Lean verification benchmarks, with thousands rather than hundreds of problems. Our work is motivated by the broader context of AI safety proposals premised on AI-driven formal verification, including the Guaranteed Safe AI agenda \cite{dalrymple2024guaranteed}.
